%! The Idea of Elimination — Unit 2, Lesson 2.1 — Homework
\documentclass[10pt]{article}
\usepackage{linalg-article}
\usepackage{linalg-boxes}
\newcommand{\TallMath}[1]{$\displaystyle #1\rule[-1.4em]{0pt}{3.2em}$}
\newcommand{\xx}{\mathbf{x}}
\newcommand{\bb}{\mathbf{b}}

\begin{document}

\pageheader{Unit 2, Lesson 2.1}{Homework}
\namedateperiod

\vspace{0.12in}

\begin{notesbox}{Practice}
Show your steps and \textbf{check} every solution by rebuilding $\bb$.

\begin{enumerate}[leftmargin=1.9em, itemsep=12pt]
  \item \textbf{Just the multiplier.} For each pair, name the \emph{pivot} (top-left coefficient) and
        the \emph{multiplier} $\ell=$ entry $\div$ pivot that would clear the term below it. Do not solve.
    \begin{enumerate}[label=(\alph*), leftmargin=1.8em, itemsep=4pt]
      \item $\begin{aligned}[t] 3x + y &= \phantom{0} \\ 6x + 5y &= \phantom{0} \end{aligned}$ \hfill
      \item $\begin{aligned}[t] 2x + 5y &= \phantom{0} \\ 8x + y &= \phantom{0} \end{aligned}$ \hfill
      \item $\begin{aligned}[t] x - 2y &= \phantom{0} \\ 5x + y &= \phantom{0} \end{aligned}$
    \end{enumerate}
    \vspace{0.9cm}
  \item \textbf{Solve by elimination.} Find $x$ and $y$ for each system.
    \begin{enumerate}[label=(\alph*), leftmargin=1.8em, itemsep=4pt]
      \item $\begin{aligned}[t] x + y &= 5 \\ 2x + 3y &= 13 \end{aligned}$ \hfill
      \item $\begin{aligned}[t] 2x + 3y &= 8 \\ 4x + 5y &= 14 \end{aligned}$ \hfill
      \item $\begin{aligned}[t] x - y &= 1 \\ 2x + y &= 8 \end{aligned}$
    \end{enumerate}
    \vspace{1.6cm}
  \item \textbf{Feed blend.} Each bag of \textbf{Feed A} has $2$ kg corn and $1$ kg oats; each bag of
        \textbf{Feed B} has $1$ kg corn and $4$ kg oats. A farmer needs exactly $13$ kg corn and
        $24$ kg oats.
    \begin{enumerate}[label=(\alph*), leftmargin=1.8em, itemsep=6pt]
      \item Set up $A\xx=\bb$ and solve by elimination for the number of bags of each feed.
            (Reorder the equations first so the multiplier is a whole number.) \vspace{1.3cm}
      \item \emph{Interpret:} how many bags of each, and does the answer make sense? \writeline
    \end{enumerate}
  \item \textbf{Three unknowns.} Solve by elimination --- clear column~1, then column~2, then
        back-substitute:
        \[
          \begin{array}{r@{}c@{}l}
            x &{}+ 2y + z ={}& 8 \\
            x &{}+ 3y + 2z ={}& 13 \\
            2x &{}+ 5y + 6z ={}& 30
          \end{array}
        \]
        \vspace{1.8cm}
  \item \textbf{Fix the zero pivot.} Elimination on $\ 0x + 4y = 8,\ \ x + 3y = 9\ $ can't start.
        Say what is wrong, fix it with a \textbf{row exchange}, and solve. \writeline
\end{enumerate}
\end{notesbox}

\begin{extensionbox}
\textbf{Meet $L$.} In Problem~4 you used three multipliers: $\ell_{21}$ and $\ell_{31}$ (Round~1)
and $\ell_{32}$ (Round~2). List all three, then drop them into a lower-triangular array with $1$s
on the diagonal:
\[
  L = \begin{bmatrix} 1 & 0 & 0 \\ \ell_{21} & 1 & 0 \\ \ell_{31} & \ell_{32} & 1 \end{bmatrix}.
\]
You have just recorded the \emph{whole} elimination in one matrix --- the $L$ of $A=LU$ coming in
Lesson~2.3. \writeline
\end{extensionbox}

\begin{spiralbox}
\textbf{Coming next --- Lesson 2.2: Elimination Matrices and Inverse Matrices.} Today each step was
``subtract $\ell$ times a row.'' Next we show that step is itself a \textbf{matrix} multiplying the
system --- and running the steps backward leads to the \textbf{inverse} of $A$.
\end{spiralbox}

\end{document}
